\subsection{FSM (Jonathan Valade)}
Afin de permettre une meilleure gestion de la FSM, les états ont été retravaillés afin de retirer le mode \textbf{Reculon} lorsque le robot perd la ligne.

Cette amélioration permet d'éviter des accélérations et désaccélérations inutiles lors de la perte de la ligne.

\subsection{Configuration variable (Tristan Faucher)}
Les modifications apportées à la bibliothèque partagée sont finalisé et un test à été effectué avec le client externe.

\subsection{Outils d'interprétation}
Afin de simplifier le post-traitement des informations, reçu par le robot,
L'équipe c'est équipé d'outils rapidement développé par l'IA afin de lire la télémétrie en temps réel et modéliser les actions prises par le robot après l'utilisation de celui-ci.

Quelques problèmes ont été noté vis-à-vis les programmes fabriqués.
Ceux-ci se sont réglé en analysant correctement la logique du code.